\chapter{Servicios de operación y niveles de servicio}
\label{cap:sd10}

Este capítulo reúne el modelo de operación de la plataforma de Distribuidora Puelche S.A. y las condiciones con las que se medirá su servicio. Articula la arquitectura del Subdocumento 4 con la dotación del Subdocumento 7, las pruebas de SD9 y la reproducción de incidentes de SD13. Su horizonte contractual comprende los 36 meses de Operación, del mes 21 al 56, y explica el enlace con el soporte puente de Etapa 1 (Distribuidora Puelche S.A., 2026a, art. 17, pp. 12--13; 2026d, cap. 10).

Se usa la definición del problema y de sus actores de SD2 y SD3 y la organización de LafroX de SD1. La responsabilidad sobre datos conserva el diseño de SD5, los cambios siguen SD6 y los riesgos mantienen su registro en SD8. E2 se consultó para cotejar capacidad de soporte. Las ramas de SD11 y SD12 consultadas aún contienen la plantilla común. Sus planes y nominaciones permanecen como dependencias documentales abiertas (LafroX, 2026a--2026l).

El documento es un borrador construido con fuentes publicadas. Los objetivos y las capacidades se identifican como compromisos o supuestos de diseño. La compilación no acredita contratación, turnos asignados, pruebas ejecutadas ni aceptación del CLIENTE. El archivo separado \textit{LAFROX-Subdocumento10-Anexos.pdf} contiene el cálculo reproducible, el contraste y las brechas. La aclaración del Informe 3 exige que la capacidad valorada se rastree hasta la propuesta técnica y conserva el horizonte de 56 meses (Moya Villegas, 2026, apartados 1 y 3).

\section{Servicios de operación}
\label{sec:sd10-servicios}
La operación convierte el diseño en un servicio atendido y medible. Cada alerta se relaciona con un proceso de negocio, un responsable y una guía ejecutable. El modelo desarrolla las prácticas de mesa de servicio, incidentes, solicitudes, problemas, cambios, conocimiento y mejora continua exigidas sobre ITIL 4. Su evidencia son tickets, bitácoras, versiones, pruebas y reportes (Distribuidora Puelche S.A., 2026b, RT-14.01--RT-14.09 y RT-21.01--RT-21.18, pp. 27 y 35--37; 2026d, apartado 10.1).
\parte{partes/10_1_servicios}

\section{Niveles de servicio}
\label{sec:sd10-niveles}
Los niveles de servicio vinculan el impacto de una falla con su atención, restitución y medición. Se mantienen los cuatro niveles del artículo 78, la clasificación funcional de SD4 y los parámetros del caso. La disponibilidad de una transacción, la severidad de un incidente y la disponibilidad de un componente se registran por separado (Distribuidora Puelche S.A., 2026a, arts. 78--79, p. 40; 2026b, RT-10.01--RT-10.02 y RT-14.03; LafroX, 2026d, apartado 4.2.4.3).
\parte{partes/10_2_niveles}
\section*{Referencias}
\addcontentsline{toc}{section}{Referencias}
Las fuentes siguientes sustentan las cifras, responsabilidades y procedimientos del borrador. La trazabilidad separada registra sus commits y localizaciones.

\begin{description}[style=nextline,leftmargin=0pt,labelwidth=0pt,labelsep=0pt,itemsep=1.5pt,parsep=0pt,topsep=4pt,font=\normalfont\intersemibold]
\item[Distribuidora Puelche S.A. (2026a).] \textit{Bases Administrativas. Licitación TFEP-01/2026}. Arts. 5, 17, 20, 21, 40, 50, 78 y 79 y Formulario T-7.
\item[Distribuidora Puelche S.A. (2026b).] \textit{Bases Técnicas Transversales}. Capítulos 7, 10, 11, 14, 19, 20, 21 y 22.
\item[Distribuidora Puelche S.A. (2026c).] \textit{Bases Técnicas del caso 02. Logística}. Capítulos 6, 10, 13, 14, 15 y 17 y Anexo B.
\item[Distribuidora Puelche S.A. (2026d).] \textit{Aclaraciones de la Licitación}. Archivos, índice, redacción, figuras, tablas, referencias y capítulo 10.
\item[Moya Villegas, A. (2026).] \textit{Informe 3. Aclaraciones Licitación TFEP-01/2026}. Comunicación del 1 de noviembre de 2026. Horizonte contractual y trazabilidad de la capacidad.
\item[LafroX (2026a).] \textit{Subdocumento 1. Presentación de la empresa}. Rama de corrección, commit \texttt{278a4c41}.
\item[LafroX (2026b).] \textit{Subdocumento 2. Problema y necesidad}. Rama de corrección, commit \texttt{a72407bc}.
\item[LafroX (2026c).] \textit{Subdocumento 3. Solución propuesta}, anexos y T-12. Rama de corrección, commit \texttt{887ab84c}.
\item[LafroX (2026d).] \textit{Subdocumento 4. Arquitectura}, anexos y T-11. Rama de corrección, commit \texttt{0dc50b3d}.
\item[LafroX (2026e).] \textit{Subdocumento 5. Modelo y gestión de datos}. Rama de corrección, commit \texttt{daec1657}.
\item[LafroX (2026f).] \textit{Subdocumento 6. Metodologías}, T-9 y T-10. Rama de corrección, commit \texttt{68a9f14f}.
\item[LafroX (2026g).] \textit{Subdocumento 7. Plan de trabajo, EDT, cronograma e implantación}, T-14, T-15 y T-18. Rama de corrección, commit \texttt{5e78b998}.
\item[LafroX (2026h).] \textit{Subdocumento 8. Plan de riesgos} y T-16. Rama de revisión, commit \texttt{1f2d739e}.
\item[LafroX (2026i).] \textit{Subdocumento 9. Plan de calidad}, T-13 y T-17. Rama de corrección, commit \texttt{fd88261e}.
\item[LafroX (2026k).] \textit{Subdocumento 13. Innovaciones}, anexos y T-19. Rama de corrección, commit \texttt{f6b57fa6}.
\item[LafroX (2026l).] \textit{Oferta Económica E2. Análisis económico-financiero}, sección 2.6. Commit \texttt{2b04aaa7}.
\end{description}

\Needspace{0.38\textheight}
\section*{Declaración de uso de IA}
\addcontentsline{toc}{section}{Declaración de uso de IA}
Se utilizó OpenAI Codex para localizar las fuentes, organizar y redactar el borrador, reproducir los cálculos y comprobar la compilación y la presentación del PDF. La contribución de generación y organización es alta. La lectura de las fuentes y las verificaciones automáticas sustentan el cotejo documental. Las brechas conservan las ausencias de evidencia.

Este documento requiere revisión humana técnica y contractual antes de su presentación final. La versión 0.1 no registra una revisión humana completada ni una aprobación del CLIENTE. Los espacios de firma permanecen sin firma. La clasificación como borrador y esta declaración se aplican también a sus anexos.

